\documentclass[letterpaper]{article} % DO NOT CHANGE THIS
\usepackage[submission]{aaai2027}  % DO NOT CHANGE THIS
% The serif, sans-serif, and monospaced fonts are loaded automatically by
% aaai2027.sty (newtxtext, helvet, courier). DO NOT add \usepackage{times},
% \usepackage{helvet}, \usepackage{courier}, or any other font package.
\usepackage[hyphens]{url}  % DO NOT CHANGE THIS
\usepackage{graphicx} % DO NOT CHANGE THIS
\urlstyle{rm} % DO NOT CHANGE THIS
\def\UrlFont{\rm}  % DO NOT CHANGE THIS
\usepackage{natbib}  % DO NOT CHANGE THIS AND DO NOT ADD ANY OPTIONS TO IT
\usepackage{caption} % DO NOT CHANGE THIS AND DO NOT ADD ANY OPTIONS TO IT
\frenchspacing  % DO NOT CHANGE THIS
\usepackage[english]{babel}

\usepackage{amsmath, amssymb, amsfonts}
\usepackage{amsthm}
\usepackage{bm}
\usepackage{booktabs}
\usepackage{array}
\usepackage{multirow}
\usepackage{tabularx}
\usepackage{mathtools}
\usepackage{microtype}
\usepackage{nicefrac}
\usepackage{makecell}
\usepackage{xcolor}
\usepackage[hidelinks]{hyperref}
\usepackage[capitalize,noabbrev]{cleveref}
\usepackage{algorithm}
\usepackage{algorithmic}
\usepackage[inline,shortlabels]{enumitem}
\usepackage{tikz}
\usetikzlibrary{arrows.meta,positioning,fit,calc,shapes.geometric}
\usepackage{pgfplots}
\pgfplotsset{compat=1.17}
\usepackage{pifont}

\pdfinfo{
/TemplateVersion (2027.1)
}

\setcounter{secnumdepth}{2}

\newcommand{\E}{\mathbb{E}}
\newcommand{\R}{\mathbb{R}}
\newcommand{\I}{\mathbb{I}}
\newcommand{\probP}{\text{I\kern-0.15em P}}

\newcommand{\cmark}{\ding{51}}%
\newcommand{\xmark}{\ding{55}}%
\let\autoref\Cref
\addto\extrasenglish{
    \def\figureautorefname{Figure}
    \def\tableautorefname{Table}
    \def\algorithmautorefname{Algorithm}
    \def\sectionautorefname{Section}
    \def\subsectionautorefname{Subsection}
    \def\proofoutlineautorefname{Proof Outline}
}

\theoremstyle{plain}
\newtheorem{theorem}{Theorem}[section]
\newtheorem{proposition}[theorem]{Proposition}
\newtheorem{lemma}[theorem]{Lemma}
\newtheorem{corollary}[theorem]{Corollary}
\theoremstyle{definition}
\newtheorem{definition}[theorem]{Definition}
\newtheorem{assumption}{Assumption}[section] % own counter: does not shift Theorem/Corollary/Proposition numbers
\theoremstyle{remark}
\newtheorem{remark}[theorem]{Remark}
\newtheorem{claim}[theorem]{Claim}
\newtheorem{hypothesis}[theorem]{Hypothesis}

\title{Discovering and Certifying Partial Symbolic Invariants \\ for Neuro-Symbolic Multi-Agent World Models}
\author{Anonymous Submission}
\affiliations{}

\begin{document}
\maketitle
\begin{abstract}
  Neuro-symbolic world models improve long-horizon reliability by injecting partial symbolic transition knowledge, such as rules for spatial consistency, object persistence, or action-conditioned effects, into an otherwise learned multi-agent World Model (MAWM). Existing approaches hand this covered/residual decomposition to the system, either hand-written or LLM-proposed and manually curated, leaving open whether it can be obtained automatically and trusted without a human in the loop. This paper studies a falsification-based discovery procedure that treats every candidate invariant as \emph{not yet refuted} rather than \emph{true}: candidates are proposed from an interaction dataset by an enumerative, tree-based, or LLM-guided proposer, certified through active and multi-policy falsification with an explicit statistical validity bound, and injected with a confidence-weighted, retractable coupling. The analysis gives six results: a survival-based validity certificate, a search-cost bound on expected false discoveries, a chaining argument to rollout fidelity, a degradation bound under policy shift, a non-identifiability theorem showing passive discovery can certify invariants that are artifacts of the data-collection policy rather than the environment, and a sufficient condition under which active, multi-policy falsification restores a finite transfer guarantee despite that impossibility. On GridCraft, Overcooked, Predator-Prey, and SMACv2, evaluated under feature anonymization and averaged over 5 seeds, a decision-tree proposer recovers $0.837\pm0.039$ of the fidelity gain of hand-written oracle rules and an LLM-guided proposer recovers $0.774\pm0.057$, both ahead of brute-force enumeration ($0.612\pm0.048$), whose larger candidate count buys no certification-budget advantage but yields more low-precision survivors. The difficulty gradient supports the non-identifiability result: a rule over a latent, never-observed environment variable is falsely certified under passive falsification alone ($\hat\varepsilon=0.006\pm0.005$) and flagged non-certifiable once active, multi-policy falsification is added, as \autoref{thm:t5} predicts.
\end{abstract}

\noindent\textbf{Keywords:} Neuro-symbolic learning, World Models, multi-agent, model-based RL, rule discovery, partial observability

\section{Introduction}
\label{sec:intro}

Model-based reinforcement learning (MBRL) supports efficient planning by learning predictive World Models (WMs)~\cite{hafner2018planet,hafner2019dreamer,janner2019mbpo,schrittwieser2020muzero}, but rollout errors compound over long horizons, especially in multi-agent settings with partial observability and combinatorial joint dynamics~\cite{willemsen2021mambpo}. A recent line of neuro-symbolic work addresses this by decomposing the observation into a \emph{symbolically covered} sub-vector, predicted by rules, and a \emph{residual} sub-vector, left to a neural component~\cite{balloch2023neurosymbolicworldmodelsadapting,zhu2022iaem}; we refer to this general recipe as \emph{covered/residual decomposition}. It reduces the hypothesis space the neural model must fit and removes error on structurally predictable transitions.

Every version of this recipe we are aware of treats the decomposition as an input: a domain expert writes the rules, or a large language model proposes them and a human filters the result~\cite{zhou2024walle,athalye2026pixelstopredicates}. This is the point on which the approach's scalability depends, and no existing method offers a guarantee for it. \emph{Can the covered/residual split be obtained automatically, and can the resulting rules be trusted without a human in the loop?}

The contribution is a procedure for \emph{discovering} the decomposition itself, together with a certificate of when a discovered invariant can be trusted. We reuse existing covered/residual injection mechanisms and focus on the discovery and certification problem.

\paragraph{Contribution.} We formalize a \emph{partial invariant} as a triplet $(\varphi, c, \rho)$: a feature reading, a precondition, and a state-and-action-indexed transport. Its validity is a conditional determinism rate rather than a binary truth value (\autoref{sec:invariants}). Invariants are proposed by an enumerative, decision-tree, or LLM-guided proposer over a compact transport catalog, then subjected to a falsification protocol of increasing severity: passive holdout, active exploration, cross-variant, and cross-generator testing. Surviving invariants are injected with a confidence interpolating between a soft regularizer and a hard constraint, and are \emph{retracted} the moment a counterexample appears during training.

We characterize this procedure with six results (\autoref{sec:theory}): a survival-based validity certificate (\autoref{thm:t1}) whose statistical cost scales with the number of candidates actually proposed rather than the size of the search space (\autoref{thm:t2}), together with a corollary showing that a proposer's expected count of false discoveries is controlled by its \emph{precision}, not by the candidate count alone (\autoref{cor:t2b}); a chaining bound from discovery validity to long-horizon rollout fidelity (\autoref{thm:t3}); a degradation bound under policy shift via a restricted concentrability coefficient (\autoref{thm:t4}); a non-identifiability theorem (\autoref{thm:t5}) showing that under partial observability, a rule can be a certified artifact of the data-collection policy rather than a fact about the environment, invisible to any passive certification procedure regardless of sample size; and a sufficient condition (\autoref{thm:t7}) under which certifying against the falsification-gate family restores a finite transfer guarantee despite that impossibility. The non-identifiability result motivates active and multi-policy falsification: two invariants with identical training-time validity can behave differently under a shifted policy.

We reuse three existing neuro-symbolic integration strategies: symbolic consistency regularization, symbolic projection, and residual symbolic dynamics. We also correct a design flaw in the standard formulation of the residual strategy: the residual head previously predicted uncovered features without access to the symbolic prediction on covered features, discarding information whenever the two families of features are correlated. We condition the residual head on the covered prediction and turn this correlation into an explicit falsification criterion instead of an unexamined assumption (\autoref{sec:injection}).

The paper is organized as follows: \autoref{sec:related} positions the work, \autoref{sec:background} defines the MAWM setting, \autoref{sec:invariants} introduces partial invariants and discovery, \autoref{sec:theory} gives the guarantees, \autoref{sec:injection} describes injection into MAWMs, \autoref{sec:exp} reports the evaluation protocol and results, and \autoref{sec:conclusion} states limitations.

\section{Related Work}
\label{sec:related}

\paragraph{World models and structured dynamics.} Latent WMs such as PlaNet and Dreamer~\cite{hafner2018planet,hafner2019dreamer,hafner2020dreamerv2,hafner2023dreamerv3} and transformer-based WMs~\cite{chen2022transdreamer,micheli2023transformers} learn all dynamics from data; structured variants add graph dynamics~\cite{zhang2021worldmodelgraph}, contrastive representations~\cite{kipf2020contrastivelearningstructuredworld}, or compositionality~\cite{zhang2024combo}, but do not separate or discover symbolically predictable structure from residual uncertainty.

\paragraph{Discovering invariants over learned or given variables.} A first family discovers \emph{which transformation} makes dynamics simple, given the variables; symbolic regression and equation discovery (e.g., SINDy-style methods) assume continuous, already-observed state variables, excluding discrete multi-agent features and preconditions. ConCerNet~\cite{zhang2023concernet} jointly learns a representation and a conserved quantity with a proven functional link. It is the closest discover-then-hard-inject precedent, but is restricted to continuous global conservation laws and single-agent settings. IAEM~\cite{zhu2022iaem} represents an action's effect as a residual with adaptive weighting where invariance fails, which is close to the covered/residual intuition but single-agent and without a statistical certificate. Flow Equivariant World Models~\cite{lillemark2026flowm} unify self- and object-motion as learned equivariant Lie-group flows; our canonical spatial-invariant example falls in this class, so any discovery claim must be positioned against it.

\paragraph{Learning symbolic action models over given predicates.} STRIPS/PDDL action-model learning gives the strongest existing guarantees, but over hand-defined predicates for full-information planning. Safe Action Model (SAM) learning extends to partial observability with a polynomial-time safety guarantee~\cite{lehiu2024safeactionmodelspo}; Conditional-SAM extends this to conditional effects, proving an exponential worst-case sample requirement~\cite{mordoch2024conditionalsam}; a probabilistic alternative infers preconditions/effects from noisy traces~\cite{lamanna2024nolam}. These assume discrete states, deterministic effects, and hand-defined predicates. From Pixels to Predicates~\cite{athalye2026pixelstopredicates} uses a pretrained VLM to propose a predicate vocabulary and select a compact subset, a propose-many/select-few schema related to our LLM proposer, but for single-agent robotic planning with full predicate-level abstraction rather than partial coverage of a latent multi-agent WM.

\paragraph{LLM-guided rule induction for world models.} WALL-E and WALL-E 2.0~\cite{zhou2024walle,zhou2025walle2} induce, update, and prune symbolic rules gradient-free by comparing explored trajectories against world-model predictions, composing the result with an LLM-based world model for Minecraft/ALFWorld-style agents. They are the closest system-level precedent to our discovery loop. They differ on the dimensions relevant to our claim: their world model is an LLM over text-described, single-agent environments, whereas ours is a latent WM under a Dec-POMDP; their rules are invoked by function call rather than injected into training; and induce/prune is heuristic, without falsification, a validity bound, or partial observability. Our non-identifiability result (\autoref{thm:t5}) applies to their setting as much as to ours and is, to our knowledge, not addressed there.

\paragraph{Evaluating world models.} WM evaluation typically emphasizes predictive accuracy or downstream reward, though recent work calls for broader diagnostics~\cite{Duan_2025_ICCV,Delvecchio2025p1157}. Multi-agent evaluation rarely measures semantic validity, such as rule violations or logical consistency, making it hard to distinguish locally-accurate from semantically-coherent models. $\mathrm{RVR}_{\mathrm{pre}}$, $\mathrm{RVR}_{\mathrm{post}}$, and the correction magnitude (\autoref{sec:injection}) instantiate this point: a raw neural MAWM can post low one-step loss while violating discovered rules on $15$ to $18\%$ of covered-feature transitions (\autoref{sec:exp-injection}), invisible to prediction loss alone.

\paragraph{Non-identifiability under partial data coverage.} \autoref{thm:t5}'s impossibility result instantiates, for invariant discovery, a coverage/concentrability phenomenon studied elsewhere under different names: offline RL requires a coverage condition for value-function guarantees to transfer to a target policy~\cite{chen2019information}, and off-policy evaluation under unobserved confounders is similarly non-identifiable without further assumptions~\cite{namkoong2020offpolicy}. \autoref{thm:t7}'s positive counterpart transposes the same style of fix to rule discovery as an operational, verifiable falsification protocol, F1 to F4, rather than an unverified data-collection assumption.

\paragraph{Gap.} No existing method jointly (i) discovers which features of a learned multi-agent latent WM are symbolically predictable, (ii) certifies discovered invariants with an explicit, falsification-based statistical guarantee rather than heuristic induce/prune, and (iii) characterizes when certification is impossible under partial observability rather than insufficient data. Covered/residual-style symbolic injection~\cite{balloch2023neurosymbolicworldmodelsadapting} and symbolic/organizational MARL biases~\cite{hasanbeig2018lcrl,soule2025moisemarl,alshiekh2018shield} supply the injection side; this paper supplies the discovery side and its guarantees.

\section{Background}
\label{sec:background}

Cooperative partially observable MARL is a Decentralized Partially Observable Markov Decision Process (Dec-POMDP)~\cite{oliehoek2016concise,Bernstein2000complexity} $\mathcal{M} = \langle S, \{A_i\}_{i=1}^n, T, R, \{\Omega_i\}_{i=1}^n, O, \gamma \rangle$. Agent $i$'s observation at time $t$ is a structured feature vector $\omega^i_t = (f^i_{t,1},\dots,f^i_{t,n_f})$; joint quantities are written $\omega^j_t=(\omega^1_t,\dots,\omega^n_t)$, $a^j_t=(a^1_t,\dots,a^n_t)$, with joint history $h^j_{t-1}$. A behavioral policy $\pi$ induces a distribution $d^\pi$ over transitions $x=(h^j_{t-1},\omega^j_t,a^j_t,\omega^j_{t+1})$; we write $x^-$ for the part preceding $\omega^j_{t+1}$.

A Multi-Agent World Model (MAWM) predicts $\mathcal{T}^j_\theta: \mathcal{H}^j\times\Omega^j\times A^j\to\mathcal{P}(\Omega^j)$, typically through a latent encoder/decoder pair with recurrent dynamics, following Dreamer- and PlaNet-style recurrent state-space models~\cite{hafner2018planet,hafner2019dreamer,ha2018worldmodels}: $z^j_t = \mathrm{Enc}^j(\omega^j_t)$, $\tilde h_t = f_\theta(\tilde h_{t-1}, z^j_t, a^j_t)$, $\hat\omega^j_{t+1} = \mathrm{Dec}^j(\hat z^j_{t+1})$ with $\hat z^j_{t+1}=g_\theta(\tilde h_t)$, and, in stochastic variants, a latent $s^j_t\sim p_\theta(s^j_t\mid s^j_{t-1},a^j_{t-1})$ with approximate inference $q_\phi(s^j_t\mid\omega^j_{\le t},a^j_{<t})$~\cite{hafner2020dreamerv2}. Such a MAWM may be centralized, factorized per agent, or relational, typically with centralized training and decentralized execution~\cite{oliehoek2016concise,marl-book,wang2022}. Compared with single-agent WMs, MAWMs face combinatorial joint-actions, non-stationarity, and inter-agent dependencies that increase the payoff of removing structurally predictable dynamics from the learned target.

In this covered/residual injection recipe, a fixed mask $M_d\in\{0,1\}^{n\times n_f}$ marks a covered feature subspace $\Omega^{d,j}$ predicted by a symbolic operator $\mathcal{T}^j_s$, with the complementary residual subspace $\Omega^{u,j}$ left to $\mathcal{T}^j_\theta$; the coverage ratio is $\kappa=\|M_d\|_0/(nn_f)$, ranging from $\kappa=0$ (a purely neural MAWM) to $\kappa=1$ (a complete symbolic simulator, used only as an oracle-style diagnostic rather than a realistic operating point). This paper's central move is to make $M_d$ (equivalently, $\kappa$) the \emph{output} of the discovery procedure defined next, rather than an input chosen by a human or lightly-audited LLM call.

\section{Partial Invariants: Definition and Discovery}
\label{sec:invariants}

\subsection{A unified object}
\label{sec:invariant-object}

\begin{definition}[Partial invariant]
  \label{def:invariant}
  A partial invariant is a triplet $r=(\varphi,c,\rho)$: a reading $\varphi$ of the joint observation, a precondition $c$ on $(h^j_{t-1},\omega^j_t,a^j_t)$, and a transport $\rho_{a^j_t}$ indexed by the joint action, such that
  \begin{equation}
    c(h^j_{t-1},\omega^j_t,a^j_t) \;\Longrightarrow\; \varphi(\omega^j_{t+1}) = \rho_{a^j_t}\big(\varphi(\omega^j_t), h^j_{t-1}\big).
  \end{equation}
  Its \emph{validity} is not a Boolean but the conditional violation rate under a data-generating policy $\pi$:
  \begin{equation}
    \label{eq:eps-def}
    \varepsilon^\pi_r \;=\; \Pr\nolimits_{d^\pi}\!\big[\varphi(\omega^j_{t+1}) \neq \rho_{a^j_t}(\varphi(\omega^j_t),h^j_{t-1}) \,\big|\, c\big].
  \end{equation}
\end{definition}
Three special cases of $\rho$ recover, as a single object, notions the literature treats separately: $\rho=\mathrm{id}$, $c=\top$ gives persistence/conservation (ConCerNet-style); $\rho$ a group action with $c=\top$ gives equivariance (Flow-Equivariant-style); $\rho$ an arbitrary bounded program with non-trivial $c$ gives a conditional rule (STRIPS/SAM-style, without discreteness). A rule is never ``true''; it is \emph{not yet falsified} at a given budget $N_c$, and $\varepsilon^\pi_r$ is relative to the policy $\pi$ that generated the certification data. \autoref{thm:t4} and \autoref{thm:t5} treat this dependence formally.

\subsection{Transport catalog and readings}
\label{sec:catalog}

To keep the search space small enough for \autoref{thm:t2} to bind usefully, we restrict $\rho$ to three parametric families (full estimators and closed-form fits in the supplementary material): \textbf{(A) bounded affine}, $\rho(\varphi,a)=\mathrm{clip}(\alpha\varphi+\delta_a,\ell,u)$ with $\alpha\in\{0,1\}$, covering persistence, constants, free motion, and clipped/cooldown dynamics as one parametric form; \textbf{(B) delayed copy}, $\rho=\psi_{t-m}$ for another reading $\psi$ and memory order $m\in\{0,\dots,m_{\max}\}$, covering pickup/drop and delayed effects; \textbf{(C) slot permutation}, $\varphi^{\sigma(i)}_{t+1}=\varphi^i_t$, covering agent-slot reassignment. Readings $\varphi$ are restricted to three levels closed under permutation of feature indices: atomic, pairwise-relational, and per-agent aggregates. This makes the anonymization experiment in \autoref{sec:exp} a test of the algorithm rather than of the code.

A full table of the catalog with closed-form estimators is in the supplementary material.

\subsection{Discovery loop: propose, falsify, promote}
\label{sec:loop}

Discovery uses four disjoint data partitions: $\mathcal{D}_{\mathrm{prop}}$ (structure search), $\mathcal{D}_{\mathrm{fit}}$ (parameter/threshold estimation), $\mathcal{D}_{\mathrm{test}}$ (falsification only, never used to select candidates), and $\mathcal{D}_{\mathrm{shift}}$ (held-out environment variants and alternative data generators). This separation is what makes \autoref{thm:t1} and \autoref{thm:t2} apply to $\mathcal{D}_{\mathrm{test}}$ without selection bias.

\begin{algorithm}[t]
  \caption{Discover, falsify, inject}
  \label{alg:main}
  \begin{algorithmic}[1]
    \STATE Screen features on $\mathcal{D}_{\mathrm{prop}}$ (non-vacuity screen, supplementary material).
    \STATE Propose $\hat{r}\in\hat{\mathcal{R}}$, $|\hat{\mathcal{R}}|=\hat M$, via \textbf{P1} (enumeration, depth $\beta$), \textbf{P2} (decision tree on $\mathcal{D}_{\mathrm{prop}}$), or \textbf{P3} (LLM over anonymized features).
    \STATE Fit transport parameters on $\mathcal{D}_{\mathrm{fit}}$.
    \FOR{F1 (passive holdout) $\to$ F2 (active exploration) $\to$ F3 (cross-variant) $\to$ F4 (cross-generator)}
    \STATE Test survivors on $\mathcal{D}_{\mathrm{test}}\cup\mathcal{D}_{\mathrm{shift}}$; compute $\hat\varepsilon_r$, $q_r$ (\autoref{thm:t1}, \autoref{thm:t2}).
    \ENDFOR
    \STATE Promote: $q_r<\tau_{\mathrm{soft}}\Rightarrow$ reject; $<\tau_{\mathrm{hard}}\Rightarrow$ soft (weighted $q_r$); else hard injection.
    \STATE Monitor training; on violation, \emph{retrograde} the rule. Re-promotion needs fresh data, capped at $R_{\max}$ attempts, folded into the budget as $\hat M_{\mathrm{eff}}\times K\times R_{\max}$.
  \end{algorithmic}
\end{algorithm}

The output of \autoref{alg:main} is the inventory $\{(\varphi,c,\rho,m,q,\mathrm{cov},N_c)\}$, the discovered mask $\hat M_d$, and the measured coverage $\hat\kappa$. These are the objects consumed by \autoref{sec:injection}. Full falsification-gate definitions, the non-vacuity screen, and complexity of each proposer are in the supplementary material.

\section{Theoretical Guarantees}
\label{sec:theory}

Let $N_c(r)=|\{x\in\mathcal{D}_{\mathrm{test}}: c(x^-)\}|$ and let $\hat\varepsilon_r$ denote the empirical violation rate on that set. All proofs are in the supplementary material.

\begin{assumption}[Certification data regularity]
  \label{ass:regularity}
  Fix $r$ independently of $\mathcal{D}_{\mathrm{test}}$, and let $(x_k)_{k=1}^{N_c(r)}$ enumerate, in temporal order, its elements satisfying $c$. One of: \textbf{(a) sub-sampled}, at most one transition per episode is retained, so $(x_k)$ are i.i.d.\ draws from $d^\pi(\cdot\mid c)$; or \textbf{(b) any-time-valid}, transitions within an episode may be arbitrarily correlated and $N_c(r)$ may itself be a stopping time, but the certificate is computed via a test-supermartingale rather than a fixed-sample tail bound~\cite{howard2020timeuniform}. We adopt (b) by default; (a) is a cheaper fallback when the supermartingale construction is not implemented.
\end{assumption}

\begin{theorem}[Certification by survival]
  \label{thm:t1}
  Under \autoref{ass:regularity}, if zero violations are observed among $N_c(r)$ transitions satisfying $c$, then with probability at least $1-\delta$, uniformly over $N_c(r)$ under regime (b), $\varepsilon^\pi_r \le \big(\ln(1/\delta) + O(\log\log N_c(r))\big)/N_c(r)$, the $\log\log$ term vanishing under regime (a).
\end{theorem}
Under (a), this reduces to the usual ``rule of three'' ($\delta=0.05\Rightarrow\varepsilon\le 3/N_c$), with a small, standard inflation under (b) for within-episode correlation and adaptive stopping; with $k\ge1$ violations, replace the bound with the Clopper-Pearson limit (a) or its any-time-valid analogue~\cite{howard2020timeuniform} (b). \autoref{thm:t2}, \autoref{cor:t2b}, and \autoref{alg:main}'s promotion rule all inherit \autoref{ass:regularity}.

Gates F1 to F4 test each candidate against data generated by a different, gate-specific policy: $\pi_1=\pi$ (passive), $\pi_2$ (F2's active, precondition-directed exploration), and $\pi_3,\pi_4$ (F3/F4's variant and alternative-generator policies). A rule that survives gate $k$ is therefore certified against $\pi_k$, \emph{not} against $\pi$.

\emph{Effective candidate budget.} $\hat M$ counts proposed \emph{rules}, but every rule is instantiated at a choice of promotion thresholds and memory order swept in practice ($\tau_{\mathrm{soft}},\tau_{\mathrm{hard}}\in$ a grid, $m\in\{0,\dots,m_{\max}\}$, precondition depth $\beta\in$ a grid), all fit or selected on the same data used to select candidates. We fold this into the budget rather than deferring it:
\begin{equation}
  \label{eq:mhat-eff}
  \hat M_{\mathrm{eff}} = \hat M \cdot |\mathrm{grid}(\tau_{\mathrm{soft}})|\cdot|\mathrm{grid}(\tau_{\mathrm{hard}})|\cdot(m_{\max}{+}1)\cdot|\mathrm{grid}(\beta)|.
\end{equation}
Because the union-bound cost below is only logarithmic in the budget, this correction is cheap: a $10$-point grid on each of four dimensions multiplies $\hat M_{\mathrm{eff}}/\hat M$ by $10^4$ at a cost of only $\ln(10^4)\approx9.2$ additional nats.

\begin{theorem}[Search cost, per gate]
  \label{thm:t2}
  Let $\hat{\mathcal{R}}$ be the $\hat M_{\mathrm{eff}}$ candidates (\autoref{eq:mhat-eff}) proposed from $\mathcal{D}_{\mathrm{prop}}\cup\mathcal{D}_{\mathrm{fit}}$ only, and let each be tested at up to $K=4$ gates on data disjoint across gates. With probability at least $1-\delta$, simultaneously for every $r\in\hat{\mathcal{R}}$ and every gate $k\in\{1,\dots,4\}$ it reaches, $\varepsilon^{\pi_k}_r \le \ln(\hat M_{\mathrm{eff}} K/\delta)/N_c^{(k)}(r)$.
\end{theorem}
The union bound over all $\hat M_{\mathrm{eff}}\times K$ (candidate, gate) pairs, whether or not a candidate reaches a given gate, licenses reusing the adaptive decision ``test at F2 only if F1 was passed'' without invalidating the certificate: no additional correction is needed for the adaptivity itself, only for the number of pairs in play. A rule promoted to hard injection thus carries $K$ separate guarantees $\varepsilon^{\pi_k}_r\le\hat\varepsilon^{(k)}_r$, not one guarantee about a single, unspecified policy. Re-promotion after retrogradation (\autoref{alg:main}) requires fresh falsification data disjoint from every prior test, keeping each attempt a fresh draw from the budget above; we additionally cap re-promotions at $R_{\max}$ per rule, folded into the budget as $\hat M_{\mathrm{eff}} \times K \times R_{\max}$.

The union-bound cost is only logarithmic in $\hat M_{\mathrm{eff}}$; the operative argument for a parsimonious proposer is instead the expected \emph{count} of false discoveries:
\begin{corollary}[Expected false discoveries and the certification-budget reading]
  \label{cor:t2b}
  Let $\mathcal{R}_{\mathrm{bad}}=\{r\in\hat{\mathcal{R}}:\varepsilon_r>\varepsilon^\star\}$ and $\mathrm{prec}_{\mathrm{prop}}=1-|\mathcal{R}_{\mathrm{bad}}|/\hat M_{\mathrm{eff}}$ the proposer's precision. Then $\E[|\{r\in\mathcal{R}_{\mathrm{bad}}:\text{survives } N_{\min}\text{ tests}\}|] \le (1-\mathrm{prec}_{\mathrm{prop}})\,\hat M_{\mathrm{eff}}\,e^{-\varepsilon^\star N_{\min}}$. Equivalently, keeping this expectation below $\delta$ requires $N_{\min}\ge\ln(\hat M_{\mathrm{eff}}/\delta)/\varepsilon^\star$, the same logarithmic dependence as \autoref{thm:t2}, read as a required sample size at fixed error rather than an error at fixed sample size.
\end{corollary}
At $\varepsilon^\star=0.01,\delta=0.05$, certifying one rule to this level costs $N_c\approx1378$ at $\hat M=48{,}213$ (P1's realized budget, \autoref{tab:discovery}) versus $\approx882$ at $\hat M=337$ (P2's), a factor of $1.56$, not $10^4$. A parsimonious proposer controls $(1-\mathrm{prec}_{\mathrm{prop}})\,\hat M_{\mathrm{eff}}$, the expected \emph{count} of bad candidates entering falsification, not the logarithmic per-candidate certification cost. The measured precision gap between P1 and P2/P3 (\autoref{tab:discovery}) is the operative evidence for a parsimonious proposer's advantage.

\begin{assumption}[Rollout support]
  \label{ass:support}
  The imagined rollout's state distribution at every horizon $h\le H$ remains within the support of the falsification distribution $d^\pi$ used to certify $\varepsilon_d$.
\end{assumption}

\begin{proposition}[Chaining to rollout fidelity, conditional]
  \label{thm:t3}
  Under hard injection with disjoint covered/residual masks and \autoref{ass:support}, if the symbolic (covered-feature) operator has violation rate $\varepsilon_d$ on covered features (diameter $B$) and the residual model has error $e_u$, the one-step error obeys $\E\|\hat\omega_{t+1}-\omega_{t+1}\|^2 \le \varepsilon_d B^2 + e_u$; under an $L$-Lipschitz rollout operator, the $H$-step accumulated error obeys $E_H \le \big(\sum_{i=0}^{H-1}L^i\big)(\varepsilon_d B^2+e_u)$.
\end{proposition}
Stated as a proposition, not a theorem, because \autoref{ass:support} is unestablished: imagined rollouts visit self-generated states, and we diagnose support empirically. On GridCraft, the median finite-difference one-step rollout Lipschitz estimate is $1.08$ (95th percentile $1.31$), and imagined states stay within the empirical-support proxy for at least $95.3\%$ of rollouts through $H=40$; we therefore treat the proposition as informative only at short/medium horizons. The $H=100$ curves in \autoref{sec:exp-gradient} should not be read as evidence for the chaining bound.

\begin{theorem}[Degradation under policy shift]
  \label{thm:t4}
  Let $C_{\pi'/\pi}(c)=\sup_{x^-\in\{c\}} d^{\pi'}(x^-)/d^\pi(x^-)$. Then $\Pr_{\pi'}[V_r\wedge c] \le C_{\pi'/\pi}(c)\cdot\Pr_\pi[V_r\wedge c]$, i.e., $\varepsilon^{\pi'}_r \le \varepsilon^\pi_r\cdot C_{\pi'/\pi}(c)\cdot \Pr_\pi[c]/\Pr_{\pi'}[c]$.
\end{theorem}
This quantifies the objection that a rule certified under one policy need not transfer to another: validity transfers according to the certification data's coverage of the precondition region under the new policy. Concretely, a rule certified at $\varepsilon^\pi_r=0.01$ transfers to $\pi'$ at $\varepsilon^{\pi'}_r\le0.02$ if $C_{\pi'/\pi}(c)=4$ and $\pi'$ visits the precondition region twice as often as $\pi$ ($\Pr_\pi[c]/\Pr_{\pi'}[c]=0.5$). These are two independent quantities. The bound offers no guarantee once $\pi'$ reaches states outside $\{c\}$'s support under $\pi$, since $C_{\pi'/\pi}(c)$ is then unbounded. Gate F4 does not estimate $C_{\pi'/\pi}(c)$ itself, since a supremum density ratio is unreliable to estimate from finite samples. It falsifies against deployment-time generators, obtaining a direct empirical estimate of $\varepsilon^{\pi'}_r$.

\begin{theorem}[Non-identifiability under partial observability]
  \label{thm:t5}
  There exist two Dec-POMDPs $\mathcal{M},\mathcal{M}'$ and a candidate rule $r$ such that $\varepsilon^\pi_r=0$ in both models under a policy $\pi$, yet $\varepsilon^{\pi'}_r>0$ in $\mathcal{M}$ under a different policy $\pi'$. No certification procedure with access only to $d^\pi$ can distinguish $\mathcal{M}$ from $\mathcal{M}'$, or bound $\varepsilon^{\pi'}_r$, for any $N_c$. \emph{(Constructive proof: a four-cell corridor where an agent's motion appears deterministic only because an exogenous, unobserved environment hazard under $\pi$ never occupies the blocking cell; see supplementary material.)} Increasing the memory order $m$ does not help when the confounding variable never enters any agent's observation, by construction rather than by omission.
\end{theorem}
\autoref{thm:t5} is the paper's central negative result, separating a \emph{deferred effect} (information present in $h_{t-m}$, fixed by increasing $m$) from a \emph{latent confounder} (information never observed, addressable only by falsifying against multiple policies and environment variants, i.e., gates F2 to F4 of \autoref{alg:main}). A discovery algorithm that only falsifies passively against one data-collection policy can certify rules that hold by accident of that policy. This is an instance, in the invariant-discovery setting, of a non-identifiability phenomenon studied in adjacent literatures: the impossibility of sample-efficient off-policy guarantees without a coverage/concentrability condition on the offline data~\cite{chen2019information}, and the impossibility of consistent policy evaluation under confounders unobserved by the acting agent~\cite{namkoong2020offpolicy}. The contribution is its instantiation for rule discovery under partial observability. \autoref{thm:t7} shows that the coverage condition identified in those literatures also motivates the F2 to F4 severity ladder.

\begin{theorem}[Sufficient condition for transfer]
  \label{thm:t7}
  Let $r$ be certified across the falsification-gate policies $\pi_1,\dots,\pi_K$ (\autoref{thm:t2}), i.e., $\varepsilon^{\pi_k}_r\le\varepsilon$ for all $k$, and let $\bar\pi$ denote their mixture. If a deployment policy $\pi'$ satisfies $C_{\pi'/\bar\pi}(c):=\sup_{x^-\in\{c\}}d^{\pi'}(x^-)/d^{\bar\pi}(x^-)<\infty$, then $\varepsilon^{\pi'}_r \le C_{\pi'/\bar\pi}(c)\cdot\max_k\varepsilon^{\pi_k}_r$.
\end{theorem}
Read together, \autoref{thm:t4} to \autoref{thm:t7} give three statements: degradation is bounded whenever the coverage ratio $C_{\pi'/\pi}(c)$ is finite (\autoref{thm:t4}), this ratio can be infinite under passive, single-policy certification (\autoref{thm:t5}), and certifying against the gate-policy family restores a finite ratio whenever the deployment policy's precondition region is covered by that family (\autoref{thm:t7}). This is the condition F2 to F4 are designed to test.

\begin{remark}[Composition]
  If $p$ invariants are injected on disjoint positions, the aggregate covered-feature violation rate is bounded by $\varepsilon_d\le\sum_{i=1}^p\varepsilon_{r_i}\Pr[c_i]$, not $\sum_i\varepsilon_{r_i}$ alone, since $\varepsilon_{r_i}$ is conditional on $c_i$ and $\Pr[V_i]=\varepsilon_{r_i}\Pr[c_i]$. At $\varepsilon=0.01$ per rule, $p=50$ rules, this gives $\varepsilon_d\le0.5\cdot\max_i\Pr[c_i]$: at a typical trigger frequency of $0.1$, $\varepsilon_d\le0.05$, an order of magnitude tighter than $\sum_i\varepsilon_{r_i}$ alone suggests. This makes $\tau_{\mathrm{hard}}$ calibrable per rule from a \emph{measured} quantity: to keep rule $i$'s contribution to a budget $\varepsilon_{\mathrm{target}}$ below $\varepsilon_{\mathrm{target}}/p$, hard-inject only if $\tau_{\mathrm{hard},i}\ge1-\varepsilon_{\mathrm{target}}/(p\cdot\widehat{\Pr}[c_i])$. Frequent rules need a stricter threshold, rare ones a looser one. This reframes the coverage-saturation pattern reported for $\kappa\to1$ in prior neuro-symbolic MAWM experiments as a consequence of uncontrolled aggregate error, not high coverage per se.
\end{remark}

\paragraph{Feature coverage vs.\ transition coverage.} $\kappa=\|M_d\|_0/(nn_f)$ counts covered positions, but a rule fires only when $c$ holds; otherwise the position reverts to the neural target. We therefore also report $\kappa_{\mathrm{eff}}=\E_{d^\pi}[\text{fraction of covered positions with }c\text{ true}]\le\hat\kappa$, the transition-level coverage used by \autoref{thm:t3}. On GridCraft, $\kappa_{\mathrm{eff}}$ is $0.284\pm0.013$ (P1), $0.337\pm0.010$ (P2), and $0.309\pm0.016$ (P3), preserving the same proposer ordering as $\hat\kappa$.

\section{Injecting Discovered Invariants into MAWMs}
\label{sec:injection}

We reuse the three neuro-symbolic integration strategies introduced above, now consuming the discovered mask $\hat M_d$ and inventory rather than a hand-specified one. \textbf{Symbolic consistency regularization} adds $\mathcal{L}_{\mathrm{symb}}=\|P_d\hat\omega^{pre,j}_{t+1}-\omega^{d,j}_{t+1}\|^2$ to the predictive loss, weighted by $\lambda q_r$ per rule. \textbf{Symbolic projection} assembles $\hat\omega^{proj,j}_{t+1}=\iota_d(\omega^{d,j}_{t+1})+\iota_u(P_u\hat\omega^{pre,j}_{t+1})$, replacing covered features at inference time. \textbf{Residual symbolic dynamics} restricts the neural target to $\Omega^{u,j}$. Regularization modifies only gradients (soft, no guarantee); projection modifies the inference-time output (hard, if the rule is correct); residual dynamics modifies the training target (hard, if the mask is correct, lowest cost at inference). All three now read $q_r$ from \autoref{alg:main} rather than treating coverage as binary.

\paragraph{Fixing the residual-conditioning gap.} The residual head $\mathcal{T}^{u,j}_\theta$ previously never observed the symbolic prediction $\omega^{d,j}_{t+1}$ it was assembled alongside, discarding information whenever covered/residual features are correlated. We condition it on the covered prediction:
\begin{equation}
  \hat\omega^{u,j}_{t+1} = \mathcal{T}^{u,j}_\theta\big(h^j_{t-1}, \omega^{u,j}_t, a^j_t,\, \omega^{d,j}_{t+1}\big).
\end{equation}
The non-vacuity screen additionally rejects a covered candidate whose removal degrades residual accuracy, and an ablation (\autoref{sec:exp}) masks correlated residual features directly rather than only discussing the effect as a limitation.

\paragraph{Rule violation rate, corrected.} We report $\mathrm{RVR}_{\mathrm{pre}}$ (raw-prediction violation rate before correction) alongside $\mathrm{RVR}_{\mathrm{post}}$ (after correction, $0$ by construction) and a \emph{correction magnitude} $\|P_d\hat\omega^{pre,j}_{t+1}-\omega^{d,j}_{t+1}\|$; reporting $\mathrm{RVR}_{\mathrm{post}}=0$ alone is a restatement of the injection mechanism, not a learned result. Under \textbf{residual} dynamics, $\omega^{pre,j}_{t+1}$ is a diagnostic-only auxiliary prediction (not the deployed output) measuring what the neural component would predict absent the rule. Thus $\mathrm{RVR}_{\mathrm{pre}}$ and the correction magnitude are identical between the Residual (fixed) and Residual (unconditioned) conditions (\autoref{sec:exp-injection}); both are computed on covered features, which the conditioning fix does not touch.

\section{Experiments}
\label{sec:exp}

\subsection{Setup and headline comparison}
\label{sec:exp-setup}
Let $B_1$ be a purely neural MAWM, $B_2$ the symbolically-injected MAWM with hand-written oracle rules, and $M$ the symbolically-injected MAWM with the mask and rules output by \autoref{alg:main}. The headline quantity is the fraction of the oracle's gain recovered by discovery:
\begin{equation}
  \text{fraction recovered} = \frac{M - B_1}{B_2 - B_1}.
\end{equation}
Discovery is evaluated under \textbf{anonymization}: dimension permutation, semantic name stripping, and injected distractor/constant features, so a proposer cannot succeed by matching familiar names. We do not apply arbitrary linear mixing of coordinates, since that would make invariants inexpressible in the catalog of \autoref{sec:catalog}; this representation dependence is reported as a limitation. On GridCraft~\cite{gridcraft2026anonymous}, invariants span a \textbf{difficulty gradient}: spatial/boundary $\to$ persistence $\to$ conservation $\to$ long-memory cooldowns $\to$ rules over a latent, never-observed environment hazard, the direct empirical target of \autoref{thm:t5} (requiring the confound outside every agent's observation, not merely one). Overcooked~\cite{Micah2020overcooked}, Predator-Prey~\cite{lowe2017multi}, and SMACv2~\cite{Ellis2023smacv2} serve as transfer environments, each with its own oracle rule set, so the comparison uses an independently-written $B_2$ per environment. All experiments use 5 seeds with the P2/residual configuration as default unless noted.

\paragraph{Matching discovered rules to the oracle.} We count a discovered rule as matching an oracle rule when predicted next-values agree on at least $99\%$ of $\mathcal{D}_{\mathrm{test}}$ transitions satisfying either precondition. Sensitivity at $95\%$/exact agreement preserves the ranking: P2 changes from $0.891/0.812$ precision/recall to $0.907/0.829$ and $0.861/0.781$, respectively.

\subsection{Discovery recovers most of the oracle's gain}
\label{sec:exp-recover}

\begin{table}[t]
  \centering\small
  \caption{Discovery vs.\ oracle rules on GridCraft, aggregated over 5 seeds (mean $\pm$ standard error).}
  \label{tab:discovery}
  \resizebox{\columnwidth}{!}{%
    \begin{tabular}{@{}lccccc@{}}
      \toprule
      Proposer                  & Fraction recov.\  & Precision       & Recall          & $\hat M$   & $\hat\kappa$    \\
      \midrule
      P1 (enumerative)          & $0.612\pm0.048$   & $0.706\pm0.031$ & $0.658\pm0.042$ & $48{,}213$ & $0.331\pm0.014$ \\
      P2 (decision tree)        & $0.837\pm0.039$   & $0.891\pm0.026$ & $0.812\pm0.033$ & $337\pm18$ & $0.392\pm0.011$ \\
      P3 (LLM-guided)           & $0.774\pm0.057$   & $0.828\pm0.045$ & $0.741\pm0.038$ & $119\pm14$ & $0.358\pm0.017$ \\
      \midrule
      Oracle ($B_2$), reference & $1.00$            & n/a             & n/a             & n/a        & $0.42$          \\
      \bottomrule
    \end{tabular}%
  }
\end{table}

\autoref{tab:discovery} shows the two parsimonious proposers recovering more of the oracle's benefit than brute-force enumeration: P2 and P3 reach $0.837\pm0.039$ and $0.774\pm0.057$ against P1's $0.612\pm0.048$. This gap is large relative to its combined standard error ($t=3.64$ for P1 vs.\ P2, $t=2.17$ for P1 vs.\ P3), the one comparison in this table that 5 seeds support beyond a suggestive level. This is the quantity \autoref{cor:t2b} identifies as operative: P1's precision ($0.706\pm0.031$) is markedly lower than P2's ($0.891\pm0.026$) despite comparable recall, so P1's much larger $\hat M$ ($48{,}213$ against $337\pm18$ and $119\pm14$) translates into a larger \emph{count} of low-precision candidates entering falsification, not a harder per-candidate certification budget, which grows only logarithmically in $\hat M$ (\autoref{thm:t2}). P2 numerically exceeds P3 here, but the gap ($t=0.91$) is not distinguishable from noise at 5 seeds; \autoref{sec:exp-significance} reports the same caveat on the transfer environments. No proposer reaches the oracle's coverage ($\hat\kappa$ between $0.331$ and $0.392$ against $\kappa=0.42$): the anonymized, factorized catalog of \autoref{sec:catalog} cannot express every oracle rule.

\subsection{Statistical significance and transfer environments}
\label{sec:exp-significance}

Treating the 5 seeds as unpaired, only P1's disadvantage against P2 and P3 clears $|t|\ge1.96$; every other comparison, including across all three transfer environments, is not distinguishable from noise ($|t|<2$ throughout). Full per-comparison statistics and a paired-seed bootstrap over the matched GridCraft runs are in the supplementary material; the bootstrap tightens the P1-vs-P2/P3 conclusion but does not change the status of P2 vs.\ P3, projection vs.\ residual, or the residual-conditioning comparison. The abstract's headline numbers ($0.837$, $0.774$, $0.612$) are GridCraft-specific; we do not generalize them to transfer environments.

\subsection{The difficulty gradient exercises \autoref{thm:t5}}
\label{sec:exp-gradient}

On GridCraft, invariant classes span a difficulty gradient, each certified under F1-only and full F1 to F4 falsification (full per-class table in the supplementary material). The planted latent-confound class is the case \autoref{thm:t5} predicts: under F1 alone, a rule over a hazard never part of any agent's observation survives falsification at $\hat\varepsilon=0.006\pm0.005$, indistinguishable from zero, since nothing in F1-only data can reveal a confound outside every observation channel; F1 to F4 flags it non-certifiable once F3-style cross-variant data exposes the rule to states where the confound is active, illustrating \autoref{thm:t7}'s coverage condition being met rather than assumed. This supports \autoref{thm:t5} and \autoref{thm:t7} on one environment and one confound structure. The long-memory cooldown class shows a different, non-confounding failure: certification fails at $m=1$ ($\hat\varepsilon=0.224\pm0.052$) and succeeds at $m=4$, distinguishing a \emph{deferred effect} from a latent confounder because the relevant information eventually enters the observation history. The remaining classes (spatial/boundary $0.011\pm0.004$, persistence $0.009\pm0.003$, conservation $0.038\pm0.011$) certify under F1 alone, unaffected by stronger gates. On the compounding-error curves (5 seeds, across-seed SD $\le0.03$ at $H=100$), the discovered-rule model reaches $0.21$ vs.\ the neural baseline's $0.52$ (vs.\ the oracle's $0.13$) at $H=100$, consistent with $\hat\kappa<\kappa$; this saturation is not evidence for \autoref{thm:t3}'s chaining bound, which is informative only at shorter horizons.

The difficulty gradient shows that the ladder changes the outcome on one planted confound; it is not a systematic evaluation of the ladder. The supplementary material therefore reports the ladder evidence separately from three broader stress tests: an F1-only vs.\ F1 to F4 ablation on \emph{fraction recovered}, F2's active-exploration interaction cost in the same units as the wall-clock comparison above, and a sweep over planted confounds of varying detectability testing \autoref{thm:t7}'s sufficient condition.

\subsection{Injection strategy comparison}
\label{sec:exp-injection}

Comparing fraction recovered across strategies (regularization $0.679\pm0.052$, projection $0.793\pm0.038$, residual fixed $0.837\pm0.039$, residual unconditioned $0.762\pm0.048$; full $\mathrm{RVR}$/correction-magnitude table in the supplementary material) isolates the residual-conditioning fix of \autoref{sec:injection}: conditioning raises fraction recovered numerically, but at $t=1.21$ this is not distinguishable from noise at 5 seeds, whereas the correlation-masking ablation (\autoref{sec:exp-ablations}) gives a cleaner effect ($-0.147\pm0.028$) for the same mechanism. Regularization trails residual dynamics ($t=-2.43$) and is borderline against projection ($t=1.77$); projection and residual are not distinguishable ($t=-0.81$). $\mathrm{RVR}_{\mathrm{post}}=0$ for both hard strategies by construction; regularization's soft constraint leaves $\mathrm{RVR}_{\mathrm{post}}=0.112\pm0.018$. We decompose P2's rule errors by replaying oracle labels over the seed-level inventory: false-positive coverage accounts for $0.071\pm0.014$ loss in fraction recovered, false-negative coverage for $0.049\pm0.012$, confirming wrongly hard-covered positions cost more than leaving a feature to the residual head.

\subsection{Ablations}
\label{sec:exp-ablations}

The false-discovery rate under F1 to F4 ($0.031\pm0.009$) is well below what \autoref{cor:t2b} projects from $\hat M$ alone at P1's scale, since $(1-\mathrm{prec}_{\mathrm{prop}})\hat M_{\mathrm{eff}}$ is small for P2/P3 due to high precision, not small $\hat M_{\mathrm{eff}}$. The retrogradation rate ($0.021\pm0.006$/1,000 steps) is low and did not vary across strategies, as expected if retractions trace mostly to small residual $\hat\varepsilon$ already visible at certification. Wall-clock favors P2 ($6.2\pm0.8$ min) over P3 ($18.4\pm2.6$) and both over P1 ($41.3\pm3.7$), so the parsimonious advantage is not bought with compute. Masking correlated residual features drops fraction recovered by $0.147\pm0.028$. Raising precondition depth $\beta{=}2{\to}3$ multiplies $\hat M$ by $6.1\times$ for a gain of only $0.024\pm0.018$, an unfavorable trade, so we keep $\beta=2$. A tabulated summary is in the supplementary material.

\paragraph{What this section does not establish.} We report no full system comparison against WALL-E-style induce/prune~\cite{zhou2024walle,zhou2025walle2} or SAM-family learning~\cite{lehiu2024safeactionmodelspo,mordoch2024conditionalsam}, which differ in setting (\autoref{sec:related}). The supplementary material instead includes the narrower ablation that fixes the proposer and replaces F1 to F4 with a WALL-E-style keep-unless-contradicted heuristic on the same fraction-recovered metric.

\subsection{Experimental details}
\label{sec:exp-details}

\paragraph{Data, code, and generated assets.} Experiments use public GridCraft~\cite{gridcraft2026anonymous}, Overcooked~\cite{Micah2020overcooked}, Predator-Prey~\cite{lowe2017multi}, and SMACv2~\cite{Ellis2023smacv2}. The anonymous artifact includes the GridCraft code archive, preprocessing/training/analysis scripts, generated transition datasets, seed-level rule inventories, and oracle rule sets; released under a research-use license.

\paragraph{Splits, hyperparameters, and compute.} Each environment uses disjoint partitions $\mathcal{D}_{\mathrm{prop}}/\mathcal{D}_{\mathrm{fit}}/\mathcal{D}_{\mathrm{test}}/\mathcal{D}_{\mathrm{shift}}$ of 60k/40k/40k/20k transitions per variant/generator, seeds $\{0,\dots,4\}$. The default MAWM is a GRU transition model (hidden 256, two-layer width-256 encoder/decoder), Adam, lr $3\cdot10^{-4}$, batch 256, sequence 32, clipping 10, 200k steps. Discovery uses $\beta=2$, $m_{\max}=4$, $\pi_{\min}=0.01$, $\delta=0.05$, $\tau_{\mathrm{soft}}\in\{0.85,0.90,0.95\}$, $\tau_{\mathrm{hard}}\in\{0.97,0.99,0.995\}$, $N_{\min}=300$, $R_{\max}=2$ (final: $0.90$/$0.99$). Runs use one NVIDIA DGX Spark (Grace Blackwell GB10, 128GB unified memory), Python 3.11/PyTorch 2.5/CUDA 12.4; full software/hardware and per-run compute tables are in the supplementary material, with P2 discovery at $6.2\pm0.8$ min on GridCraft.

\paragraph{Contamination.} GridCraft, Overcooked, and SMACv2 are public, so P3's LLM may have seen their structure; two mitigations (an unpublished-environment test, value-range rescaling) are discussed in the supplementary material.

\section{Conclusion and Limitations}
\label{sec:conclusion}

This paper reframes neuro-symbolic multi-agent world modeling from injecting hand-specified symbolic knowledge to discovering and certifying it. On GridCraft, a parsimonious proposer recovers $0.837\pm0.039$ of the fidelity gain of hand-written oracle rules under full anonymization, consistent with \autoref{thm:t5}'s central negative result. The approach has limitations: the transport catalog does not express every oracle rule; \autoref{thm:t3}'s chaining argument assumes rollout-support we could not establish as a theorem; \autoref{thm:t4} bounds, but does not eliminate, degradation under policy shift; and a silently-confounded rule is indistinguishable from a genuine invariant without the F2--F4 family; code and data to recompute these results are in the supplementary material.

\bibliography{references}

\end{document}
